\documentclass[11pt,a4paper]{article}
\usepackage[margin=0.8in]{geometry}
\usepackage{booktabs}
\usepackage{hyperref}
\hypersetup{colorlinks=true,urlcolor=blue,linkcolor=black}
\title{CDP NLP Clustering: Empirical Comparison and Conditional Method Choice}
\author{}
\date{29 September 2026}
\begin{document}
\maketitle

\section{Decision}
Choose \textbf{multilingual MiniLM embeddings with k-means as the provisional
main method for organizing eligible responses for human review}.
Retain TF--IDF/NMF as the lexical robustness comparator and the density
method for discovering coherent subsets with explicit abstention.
This is a choice conditional on wanting broad response coverage and
more than a few coarse groups, not evidence that k-means has the highest
human-label accuracy. The 2024 results show only moderate initialization
stability and weak separation. None of the tested systems is ready to
replace the proposed multi-label human codebook.

The verification evaluates clustering methods on the existing multilingual
MiniLM embeddings. It does \emph{not} establish the best current embedding
model or state-of-the-art performance. Multilingual E5, BGE-M3, and Qwen3
Embedding are credible newer encoder candidates, but were not run here.

\section{What was actually run}
The new script is \path{src\cdp_text_clustering\verify_cdp_clustering_choice.py}.
It runs three methods on six datasets at five initialization seeds
(42--46): \textbf{90 fits}, plus 18 single-seed k-means granularity checks.
The three methods are TF--IDF/NMF, direct cached MiniLM/k-means, and
MiniLM/UMAP/HDBSCAN. The last reproduces BERTopic's \emph{clustering core},
not its topic-word representation or a newly executed full BERTopic package.
The earlier repository benchmark contains actual BERTopic runs.

For each dataset, companies connected through identical normalized text
are grouped together. The script deduplicates normalized text, selects
at most 1,800 responses, then creates a fixed approximately 80/20
group holdout. No company or normalized-exact-text group crosses
training and holdout \emph{within a dataset}. Separate datasets are not
independent replications: in particular, their 2024 initiative populations
overlap. Near duplicates and corporate-parent relationships are not resolved.
The holdout is reused across seeds and granularity checks, so this is
development-stage internal validation, not an untouched final test.

\begin{center}
\small
\begin{tabular}{lrrr}
\toprule
Dataset & Available cached rows & Training & Holdout \\
\midrule
2016--2024 initiatives & 1,800 & 1,440 & 359 \\
2016--2024 risks & 1,800 & 1,437 & 363 \\
2016--2024 opportunities & 1,800 & 1,458 & 342 \\
2024 Q7.55.2 comments & 10,545 & 1,440 & 360 \\
2024 Q7.55.3 comments & 11,695 & 1,444 & 356 \\
2024 Q7.55.4 explanations & 580 & 463 & 116 \\
\bottomrule
\end{tabular}
\end{center}
One extra normalized duplicate is removed from the longitudinal initiative
sample and from Q7.55.4. Longitudinal texts include structured questionnaire
options; 2024 inputs are narratives only. Their score differences cannot
isolate the effect of removing structured options because populations and
embedding settings also differ. The longitudinal encoder cache records
256 tokens; the 2024 script did not explicitly record a sequence limit.

K-means and NMF use 12 groups, except six for Q7.55.4.
UMAP uses 15 neighbors, five dimensions, cosine distance, and zero
minimum distance. HDBSCAN uses minimum cluster size 25 (15 for Q7.55.4)
and minimum samples one. These are a bounded comparison, not an exhaustive
hyperparameter search. K-means uses ten initializations per seed.
NMF uses NNDSVD with small random initialization perturbations; high
stability under these perturbations does not imply resampling stability.

\section{New comparison results}
The following are means over five seeds. Semantic silhouette is evaluated
in the shared MiniLM space; lexical silhouette uses training-fitted
TF--IDF. Each space favors methods operating in that space.
Density silhouettes exclude unassigned responses.

\begin{center}
\small
\begin{tabular}{llrrrr}
\toprule
Dataset & Method & Outliers (\%) & Sem. sil. & Lex. sil. & Seed ARI \\
\midrule
Long. initiatives & NMF & 0.0 & 0.078 & 0.062 & 1.000 \\
 & k-means & 0.0 & 0.131 & 0.059 & 0.782 \\
 & Density & 19.4 & 0.142 & 0.073 & 0.989 \\
\addlinespace
Long. risks & NMF & 0.0 & -0.008 & 0.038 & 1.000 \\
 & k-means & 0.0 & 0.073 & 0.017 & 0.668 \\
 & Density & 0.0 & 0.169 & 0.028 & 0.996 \\
\addlinespace
Long. opportunities & NMF & 0.0 & 0.014 & 0.027 & 0.807 \\
 & k-means & 0.0 & 0.055 & 0.008 & 0.601 \\
 & Density & 38.9 & 0.065 & 0.026 & 0.464 \\
\addlinespace
Q7.55.2 & NMF & 5.6 & -0.062 & 0.016 & 0.854 \\
 & k-means & 0.0 & 0.060 & 0.009 & 0.600 \\
 & Density & 44.8 & 0.027 & -0.002 & 0.925 \\
\addlinespace
Q7.55.3 & NMF & 5.6 & -0.043 & 0.027 & 0.782 \\
 & k-means & 0.0 & 0.044 & 0.002 & 0.516 \\
 & Density & 54.9 & 0.044 & 0.008 & 0.914 \\
\addlinespace
Q7.55.4 & NMF & 6.9 & -0.071 & 0.060 & 1.000 \\
 & k-means & 0.0 & 0.047 & -0.006 & 0.544 \\
 & Density & 81.2 & 0.037 & 0.005 & 0.695 \\
\bottomrule
\end{tabular}
\end{center}

\paragraph{How to read stability.}
ARI is adjusted Rand index between assignments to the same held-out
responses across different seeds; one means identical partitions.
For methods with abstention, the displayed ARI uses responses assigned
by \emph{both} runs. The mean common-assigned fractions for density are
73.4\% (longitudinal initiatives), 100\% (risks), 45.4\% (opportunities),
45.1\% (Q7.55.2), 34.9\% (Q7.55.3), and only 11.6\% (Q7.55.4).
High conditional stability therefore does not imply whole-population
stability. All seed-pair metrics and noise-inclusive ARI are saved.
The ten pairwise comparisons are dependent and are not confidence intervals.

\paragraph{Why not simply choose the highest score?}
The density model gives longitudinal risks just \emph{two groups in every
run}. Its excellent stability and silhouette describe a coarse partition,
not a detailed 12-category taxonomy. For longitudinal opportunities,
density discovers between two and twenty groups across seeds.
For 2024 narratives, it rejects a large majority of Q7.55.3 and Q7.55.4
holdout responses. Density remains useful for coherent subsets.

NMF is more lexically separated and often more initialization-stable.
Its English word analyzer is disadvantaged on multilingual free text:
5.6--6.9\% of the 2024 holdout responses have effectively zero topic weights.
These are explicitly unassigned, rather than silently mapped to topic zero.
Negative semantic silhouettes do not alone prove incorrect labels.

K-means supplies broad coverage and finer resolution, with better semantic
separation than NMF on all six datasets, but that comparison is partly
representation-dependent. It also \emph{forces} a nearest-centroid label.
Neither zero outliers nor moderate ARI proves substantive correctness.
If interpretability or reliable abstention is the primary objective,
the preferred method may change.

\paragraph{Number of groups is not validated.}
The eight-group k-means sensitivity runs have higher semantic silhouette
than the twelve-group runs on all five applicable datasets. For Q7.55.2,
the seed-42 values are 0.063 at eight groups, 0.046 at twelve, and 0.049
at sixteen. Twelve groups are retained as the existing fine-grained
reference specification, not declared optimal. Fewer groups can improve
separation merely by merging distinct themes. Human granularity judgments
and resampling stability are needed before freezing the number of codes.

\section{Selected reference outputs: 2024 initiatives}
The seed-42 training-fitted k-means model was used to assign all 10,545
eligible cached Q7.55.2 responses directly from their embeddings.
It was \emph{not} refitted on holdout rows, and it does not assign the
7,852 initiative rows excluded by the original 80-character threshold.
The counts below describe that eligible population, not all companies
or all CDP initiatives. Topic IDs are new and do not match the prior
BERTopic or longitudinal hybrid topic IDs.

\begin{center}
\small
\begin{tabular}{rlr}
\toprule
ID & Draft English interpretation & Responses \\
\midrule
0 & Air and cooling system efficiency & 985 \\
1 & Renewable electricity procurement & 944 \\
2 & Quantified savings and investment narratives & 996 \\
3 & Heat, steam, and fuel efficiency & 720 \\
4 & General operational emissions-reduction measures & 813 \\
5 & LED and lighting upgrades & 647 \\
6 & Fleet electrification and vehicle efficiency & 683 \\
7 & Travel, workplaces, and organizational changes & 1,010 \\
8 & Solar generation projects & 890 \\
9 & Waste, materials, and circularity & 685 \\
10 & Energy-management and efficiency programs & 1,050 \\
11 & Broad carbon-reduction and savings narratives & 1,122 \\
\midrule
 & Total & 10,545 \\
\bottomrule
\end{tabular}
\end{center}

These English labels are AI drafts based on automatic words and selected
source examples, not independent human validation. Groups 2, 4, 7, 10,
and 11 particularly need splitting, merging, or multi-label coding.
Some distinguish reporting style or project portfolios rather than a
single mechanism. A solar representative describes a planned project,
illustrating that topic membership does not establish implementation status.
Examples and source text remain available for checking every label.

The no-active-initiative explanations remain unsuitable for automatic
substantive classification without review: example clusters still show
language-related patterns and multiple barriers in one answer.
Human coding of all 580 eligible explanations remains the recommendation.

\section{Modern encoders versus verified local results}
\begin{center}
\small
\begin{tabular}{lp{8.8cm}}
\toprule
Encoder & Relevant primary-source information \\
\midrule
Multilingual E5-base & 768 dimensions, 512-token input example; its authors
recommend the \texttt{query:} prefix for clustering features. \\
BGE-M3 & 1024-dimensional dense representation, over 100 languages,
up to 8192 tokens; much of its published evidence concerns retrieval. \\
Qwen3-Embedding-0.6B & 100+ languages, 32k context, up to 1024 dimensions;
its model card includes clustering among supported tasks. \\
\bottomrule
\end{tabular}
\end{center}
These model-card capabilities motivate a future paired encoder comparison;
they are not CDP clustering results. The Qwen card's leaderboard claim
refers to the 8B model and June 2025, not the 0.6B model's current rank.
No live leaderboard rank was verified. The current experiment cannot
honestly establish a 2026 state-of-the-art winner.

Primary sources:
\url{https://huggingface.co/intfloat/multilingual-e5-base},
\url{https://huggingface.co/BAAI/bge-m3}, and
\url{https://huggingface.co/Qwen/Qwen3-Embedding-0.6B}.

\section{Artifacts, reproducibility, and remaining verification}
All new results are under
\path{data/outputs/cdp_text_clustering/cdp_clustering_verification_20260929/}.
\texttt{comparison\_summary.csv} contains the comparison;
\texttt{metrics\_by\_seed.csv} and
\texttt{stability\_seed\_pairs.csv} contain the underlying diagnostics.
\texttt{split\_manifest.csv}, \texttt{dataset\_audit.csv}, and
\texttt{metadata.json} record splits, source checksums, and package versions.
The \texttt{q7\_55\_2\_draft\_english\_labels.csv} file keeps proposed
English interpretations separate from machine assignments.

Each dataset has seed-42 k-means assignments, a serialized model, summary,
and centroid-nearest plus random training examples.
The longitudinal exports contain only the 1,800 cached responses per type,
\emph{not} a refit or replacement of the 314,298-row hybrid corpus.
The 2024 exports contain all 10,545, 11,695, and 580 eligible cached rows.
Original assignment files and the user's pilot annotations remain unchanged.
The pilot's human annotations were not used for fitting or accuracy claims;
their original source texts may occur in the unsupervised sample.

The saved predictions were checked against reloaded models; source IDs,
full narrative text, cluster totals, all 90 seed/method records, and
company/exact-text split separation were verified.
Package versions are NumPy 2.4.6, pandas 3.0.5, scikit-learn 1.6.1,
UMAP 0.5.7, HDBSCAN 0.8.40, and SciPy 1.17.1.
These differ from the earlier benchmark environment; do not interpret
changes from that benchmark as a controlled estimate of the grouping effect.
The script refuses to overwrite an existing nonempty output directory;
pass a new \texttt{--output} directory for reruns.

The next scientific validation is independent, blinded, multi-label human
coding, including non-English and uncertain cases, followed by
company-resampling and temporal-transfer checks. Until then, the method
choice is supported for \emph{exploratory organization}, not verified
measurement of actual abatement or a universally best NLP classifier.
\end{document}
